\documentclass[10pt, aspectratio=169]{beamer}
\usefonttheme{professionalfonts}

\mode<presentation>
{
  \usetheme{Berkeley}
  \usecolortheme{beaver}
  \usefonttheme{default}
  \setbeamertemplate{navigation symbols}{}
  \setbeamertemplate{caption}[numbered]
} 

\setbeamertemplate{footline}{%
  \leavevmode%
  \hbox{%
    \begin{beamercolorbox}[wd=.85\paperwidth,ht=2.5ex,dp=1ex,left]{author in head/foot}%
      \usebeamerfont{author in head/foot}Maxx Seminario, Electronic Circuits, Spring 2026%
    \end{beamercolorbox}%
    \begin{beamercolorbox}[wd=.15\paperwidth,ht=2.5ex,dp=1ex,right]{date in head/foot}%
      \hspace*{0.5em}\insertframenumber{} / \inserttotalframenumber\hspace*{0.5em}%
    \end{beamercolorbox}%
  }%
  \vskip0pt%
}

\usepackage[english]{babel}
\usepackage[utf8x]{inputenc}
\usepackage{tikz}
\usetikzlibrary{shapes.geometric}
\usepackage{pgfplots}
\usepackage{array}
\usepackage{makecell}
\usepackage{verbatim}
\usepackage{graphicx}
\usepackage{subcaption}
\usepackage{amsfonts}
\usepackage{amsmath}
\usepackage{bm}
\usepackage{epstopdf}
\usepackage{circuitikz}
\usepackage{caption}
\usepackage{multirow}
\captionsetup{compatibility=false}
\usepackage[absolute,overlay]{textpos}
\usetikzlibrary{calc}
\usetikzlibrary{pgfplots.fillbetween, backgrounds}
\usetikzlibrary{positioning}
\usetikzlibrary{pgfplots.groupplots}
\usetikzlibrary{plotmarks}
\usetikzlibrary{calc}

\pgfplotsset{compat=1.16}

% Added by Maxx Seminario 01/06/2026 - for colored icons in itemize labels
\usepackage{wasysym} % for smiles and frowns
\newcommand{\neutralface}{%
  \tikz[baseline=-0.6ex]{
    \draw (0,0) circle (0.9ex);
    \fill (-0.35ex,0.25ex) circle (0.12ex);
    \fill ( 0.35ex,0.25ex) circle (0.12ex);
    \draw (-0.35ex,-0.25ex) -- (0.35ex,-0.25ex);
  }%
}

\newcommand{\baditem}{\textcolor{red!70!black}{\frownie}}
\newcommand{\gooditem}{\textcolor{green!60!black}{\smiley}}
\newcommand{\mehitem}{\textcolor{orange!80!black}{\neutralface}}


% =========================
% Solution toggle 
% =========================
\newif\ifshowsolutions
\showsolutionstrue   %  compile WITH solutions
%\showsolutionsfalse %  compile WITHOUT solutions



\title{MOSFET Amplifier Configurations}
\subtitle{Unit 5: Field-Effect Transistors}
\author{Maxx Seminario}
\institute{University of Nebraska-Lincoln}
\date{Spring 2026}

\begin{document}

\begin{frame}
  \titlepage
\end{frame}

% Slide 2: Lecture Overview
\begin{frame}{Lecture Overview}
\begin{columns}[T]
\begin{column}{0.45\textwidth}
\textbf{Review from Last Lecture}
\begin{itemize}
\item MOSFET small-signal models
\item Transconductance ($g_m$)
\item Output resistance ($r_o$)
\item AC analysis techniques
\end{itemize}
\end{column}

\begin{column}{0.45\textwidth}
\textbf{Today's Topics}
\begin{itemize}
\item Common-Source (CS) Amplifier
\item Common-Drain (CD) Amplifier
\item Common-Gate (CG) Amplifier
\item Comparison of configurations
\end{itemize}
\end{column}
\end{columns}
\end{frame}

%==============================================================================
% COMMON-SOURCE AMPLIFIER
%==============================================================================

\begin{frame}{Common-Source (CS) Amplifier}
\begin{columns}[c]
\begin{column}{0.45\textwidth}
\textbf{Basic Configuration}
\begin{center}
\begin{circuitikz}[scale=0.8, transform shape, american, arrowmos]
  \node[nmos] (M1) at (3,2) {};
  \draw (3,0) node[ground] {};
  \draw (3,4.5) node[vcc] {};
  \node at (3,5) {$V_{DD}$};
  \node[left] at (M1.gate) {$v_{in}$};
  \draw (M1.source) -- (3,0);
  \draw (3,4.5) to[R, l=$R_D$] (M1.drain);
  \draw (M1.drain) -- (3,3);
  \draw (3,3) -- (4,3);
  \draw (4,3) to[short, -o] (5,3) node[right] {$v_{out}$};
\end{circuitikz}
\end{center}
\end{column}

\begin{column}{0.50\textwidth}
\textbf{Key Characteristics}
\begin{itemize}
\item Input at gate, output at drain
\item Source is AC ground (common)
\item Most commonly used configuration
\item Provides voltage gain with inversion
\end{itemize}

\vspace{0.3cm}
\textbf{Basic Formulas}
\begin{align*}
A_v &\approx -g_m R_D \\
R_{in} &\approx \infty \\
R_{out} &\approx R_D
\end{align*}
\end{column}
\end{columns}
\end{frame}

\begin{frame}{Common-Source Amplifier: Voltage Gain}
\textbf{Voltage Gain Derivation}
\begin{itemize}
\item Output voltage: $v_{out} = -g_m v_{gs} \cdot R_D$
\item Input voltage: $v_{in} = v_{gs}$ (infinite gate impedance)
\item Voltage gain: $A_v = \frac{v_{out}}{v_{in}} = -g_m R_D$
\end{itemize}

\vspace{0.5cm}
\textbf{Key Points}
\begin{itemize}
\item \baditem Negative sign indicates 180° phase inversion
\item \gooditem Gain magnitude can be large 
\item \gooditem High input impedance (gate draws no current)
\item \mehitem Moderate output impedance 
\end{itemize}
\end{frame}

%==============================================================================
% COMMON-DRAIN AMPLIFIER (SOURCE FOLLOWER)
%==============================================================================

\begin{frame}{Common-Drain (CD) Amplifier - Source Follower}
\begin{columns}[c]
\begin{column}{0.45\textwidth}
\textbf{Basic Configuration}
\begin{center}
\begin{circuitikz}[scale=0.8, transform shape, american, arrowmos]
  \node[nmos] (M1) at (3,2) {};
  \draw (3,0) node[ground] {};
  \draw (3,4.5) node[vcc] {};
  \node at (3,5) {$V_{DD}$};
  \node[left] at (M1.gate) {$v_{in}$};
  \draw (M1.drain) -- (3,4.5);
  \draw (M1.source) to[R, l=$R_S$] (3,0);
  \draw (M1.source) -- (4,1.2);
  \draw (4,1.2) to[short, -o] (5,1.2) node[right] {$v_{out}$};
\end{circuitikz}
\end{center}
\end{column}

\begin{column}{0.50\textwidth}
\textbf{Key Characteristics}
\begin{itemize}
\item Input at gate, output at source
\item Drain is AC ground (common)
\item Also called ``source follower''
\item Output follows input (no inversion)
\end{itemize}

\vspace{0.3cm}
\textbf{Basic Formulas}
\begin{align*}
A_v &\approx \frac{g_m R_S}{1 + g_m R_S} < 1 \\
R_{in} &\approx \infty \\
R_{out} &\approx \frac{1}{g_m}
\end{align*}
\end{column}
\end{columns}
\end{frame}

\begin{frame}{Common-Drain Amplifier: Voltage Gain}
\textbf{Voltage Gain Derivation}
\begin{itemize}
\item Output voltage: $v_{out} = g_m v_{gs} \cdot R_S$
\item But $v_{gs} = v_{in} - v_{out}$ (negative feedback)
\item Solving: $A_v = \frac{v_{out}}{v_{in}} = \frac{g_m R_S}{1 + g_m R_S}$
\end{itemize}

\vspace{0.5cm}
\textbf{Key Points}
\begin{itemize}
\item \gooditem No phase inversion (output in phase with input)
\item \baditem Voltage gain always less than 1 
\item \gooditem Very high input impedance
\item \gooditem Very low output impedance ($\approx 1/g_m$, tens of $\Omega$)
\item \gooditem Excellent buffer for driving low-impedance loads
\end{itemize}
\end{frame}

%==============================================================================
% COMMON-GATE AMPLIFIER
%==============================================================================

\begin{frame}{Common-Gate (CG) Amplifier}
\begin{columns}[c]
\begin{column}{0.45\textwidth}
\textbf{Basic Configuration}
\begin{center}
\begin{circuitikz}[scale=0.8, transform shape, american, arrowmos]
  \node[nmos] (M1) at (3,2) {};
  \draw (3,0) node[ground] {};
  \draw (3,4.5) node[vcc] {};
  \node at (3,5) {$V_{DD}$};
  \draw (M1.gate) -- (2,2.5);
  \draw (2,2.5) node[left] {$V_{bias}$};
  \draw (M1.source) to[sV, l=$v_{in}$] (3,0);
  \draw (3,4.5) to[R, l=$R_D$] (M1.drain);
  \draw (M1.drain) -- (3,3);
  \draw (3,3) -- (4,3);
  \draw (4,3) to[short, -o] (5,3) node[right] {$v_{out}$};
\end{circuitikz}
\end{center}
\end{column}

\begin{column}{0.50\textwidth}
\textbf{Key Characteristics}
\begin{itemize}
\item Input at source, output at drain
\item Gate is AC ground (common)
\item Less commonly used
\item No phase inversion
\end{itemize}

\vspace{0.3cm}
\textbf{Basic Formulas}
\begin{align*}
A_v &\approx g_m R_D \\
R_{in} &\approx \frac{1}{g_m} \\
R_{out} &\approx R_D
\end{align*}
\end{column}
\end{columns}
\end{frame}

\begin{frame}{Common-Gate Amplifier: Characteristics}
\textbf{Voltage Gain Derivation}
\begin{itemize}
\item Output voltage: $v_{out} = g_m v_{gs} \cdot R_D$
\item Input voltage: $v_{in} = -v_{gs}$ (gate at AC ground)
\item Voltage gain: $A_v = \frac{v_{out}}{v_{in}} = g_m R_D$
\end{itemize}

\vspace{0.5cm}
\textbf{Key Points}
\begin{itemize}
\item \gooditem No phase inversion (positive gain)
\item \gooditem Voltage gain similar to CS amplifier
\item \baditem Low input impedance ($\approx 1/g_m$, tens of $\Omega$)
\item \mehitem Moderate output impedance ($R_D$)
\end{itemize}
\end{frame}

%==============================================================================
% COMPARISON TABLE
%==============================================================================

\begin{frame}{MOSFET Amplifier Configurations: Summary Comparison}
\begin{table}
\centering
\small
\begin{tabular}{|l|c|c|c|}
\hline
\textbf{Parameter} & \textbf{Common-Source} & \textbf{Common-Drain} & \textbf{Common-Gate} \\
\hline
\hline
\textbf{Input Terminal} & Gate & Gate & Source \\
\hline
\textbf{Output Terminal} & Drain & Source & Drain \\
\hline
\textbf{Common Terminal} & Source & Drain & Gate \\
\hline
\hline
\textbf{Voltage Gain $A_v$} & $-g_m R_D$ & $\dfrac{g_m R_S}{1 + g_m R_S} < 1$ & $g_m R_D$ \\
& (10-100) & (0.8-0.99) & (10-100) \\
\hline
\textbf{Phase Shift} & 180° (inverted) & 0° (in phase) & 0° (in phase) \\
\hline
\hline
\textbf{Input Impedance} & Very High & Very High & Low \\
$R_{in}$ & $\approx \infty$ & $\approx \infty$ & $\approx 1/g_m$ \\
\hline
\textbf{Output Impedance} & Moderate & Very Low & Moderate \\
$R_{out}$ & $R_D \| r_o$ & $\dfrac{1}{g_m} \| R_S \| r_o$ & $R_D \| r_o$ \\
\hline
\hline
\textbf{Primary Application} & Voltage amplifier & Buffer/Driver & Current buffer \\

\hline
\end{tabular}
\end{table}
\end{frame}

\begin{frame}{Choosing the Right Configuration}
\textbf{Common-Source (CS)}
\begin{itemize}
\item Use when: High voltage gain needed, phase inversion acceptable
\item Advantages: Large gain, high input impedance
\item Disadvantages: Phase inversion, moderate output impedance
\end{itemize}

\vspace{0.3cm}
\textbf{Common-Drain (CD) - Source Follower}
\begin{itemize}
\item Use when: Need to drive low-impedance loads, buffering required
\item Advantages: Unity gain, very low output impedance, no inversion
\item Disadvantages: Gain less than 1, no voltage amplification
\end{itemize}

\vspace{0.3cm}
\textbf{Common-Gate (CG)}
\begin{itemize}
\item Use when: Wideband operation needed, current buffering
\item Advantages: Good gain, no inversion, low Miller capacitance
\item Disadvantages: Low input impedance (hard to drive)
\end{itemize}
\end{frame}

\begin{frame}{Multi-Stage Amplifiers}
\textbf{Cascading Configurations}
\begin{itemize}
\item Often multiple stages are cascaded to achieve desired performance
\item Example: CS followed by CD
  \begin{itemize}
  \item CS stage provides voltage gain
  \item CD stage provides buffering and low output impedance
  \end{itemize}
\item Total voltage gain: $A_{v,total} = A_{v,CS} \times A_{v,CD}$
\item Design considerations:
  \begin{itemize}
  \item Input impedance determined by first stage
  \item Output impedance determined by last stage
  \item Overall gain is product of individual gains
  \end{itemize}
\end{itemize}
\end{frame}

\begin{frame}{Key Takeaways}
\begin{enumerate}
\item \textbf{Three basic MOSFET amplifier configurations:}
   \begin{itemize}
   \item Common-Source: High gain, inverted output
   \item Common-Drain: Unity gain buffer, very low $R_{out}$
   \item Common-Gate: High gain, low $R_{in}$, no inversion
   \end{itemize}

\vspace{0.3cm}
\item \textbf{Key trade-offs:}
   \begin{itemize}
   \item Voltage gain vs. input/output impedance
   \item Phase relationship (inverted vs. non-inverted)
   \item Application-specific requirements
   \end{itemize}

\vspace{0.3cm}
\item \textbf{Configuration selection based on:}
   \begin{itemize}
   \item Required voltage gain
   \item Input/output impedance matching
   \item Load driving capability
   \item Frequency response requirements
   \end{itemize}
\end{enumerate}
\end{frame}


\end{document}
